% Section 1 — Introduction
% Status: v1 draft. Anchor on Wearables AI CFP topic-area alignment
% and the three-paragraph contribution structure:
%   (1) The problem: streaming egocentric retrieval is unbounded;
%       brute-force CLIP grows linearly in N.
%   (2) The architecture: PSM as a spatial-temporal index +
%       reservoir-sampled exemplars + ring-buffer time decay, with
%       per_cell_cap as the tuneable place-vs-similarity knob.
%   (3) The contribution: bounded PSM approaches brute-force while the
%       full-reservoir row checks the expected cap=K equivalence.

\section{Introduction}
\label{sec:intro}

For an intelligent agent to hold persistent memory of its environment, its internal state cannot grow unboundedly with time. Standard retrieval-augmented systems embed every observation into a dense vector bank that grows linearly with recording duration, making continuous multi-session memory prohibitively expensive on-device. Long-context modeling of the models themselves~\cite{arora2024linear} does not address the systems bottleneck: how to structure a dynamic spatial representation supporting test-time memory updating without unbounded state explosion?

\paragraph{The memory-centric systems question.}
Rather than claim a new memory bound, we ask an empirical allocation question:
\emph{under a fixed global exemplar budget $M$, how should a streaming agent
decide which observations to keep, and does allocating that budget across space
(rather than across time) preserve more of the experiences a later query will
need?} This isolates \emph{allocation} from encoder quality and from total
memory: every policy we compare keeps the same number of exemplars and returns
the same number of candidates, so differences reflect what each policy chooses
to retain. \psm{} is one point in this space; its per-cell state is bounded and
grows sublinearly in the frame count, but the number of occupied cells grows
with area, so \psm{} is budgeted per cell and is \emph{not} globally
memory-bounded (\S\ref{sec:method}).

% F1 — PSM architecture overview (two-row layout for LNCS single column).
% Top row: streaming input + PSM ingest (Stages 0–1).
% Bottom row: query retrieval + MLLM rerank (Stages 2–3).
% Pure TikZ; requires \usetikzlibrary{positioning,arrows.meta,fit,
% calc,shapes.geometric,backgrounds} loaded in main.tex.
\begin{figure}[t]
\centering
\begin{tikzpicture}[
    >={Stealth[length=1.6mm,width=1.3mm]},
    font=\scriptsize,
    node distance=1.5mm and 2mm,
    box/.style={draw, rounded corners=1pt, align=center,
                inner sep=1.5pt, minimum height=4.5mm, fill=white},
    iobox/.style={box, fill=gray!8},
    stage/.style={draw, dashed, rounded corners=2pt,
                  inner sep=2.5pt, gray!70},
    stagelbl/.style={font=\scriptsize\bfseries, gray!60!black},
    hex/.style={draw, regular polygon, regular polygon sides=6,
                minimum size=5mm, inner sep=0pt, rotate=30,
                fill=gray!5},
    hexhi/.style={hex, fill=blue!12, draw=blue!60!black, thick},
    hllbox/.style={draw, minimum width=1.8mm, minimum height=1.8mm,
                   inner sep=0pt, fill=orange!25, draw=orange!70!black},
    hllold/.style={hllbox, fill=orange!8, draw=orange!30},
    resvslot/.style={draw, minimum width=1.8mm, minimum height=1.8mm,
                     inner sep=0pt, fill=blue!15, draw=blue!60!black},
    flow/.style={->, semithick, gray!60!black},
    smallflow/.style={->, thin, gray!60!black},
]

% ===================== TOP ROW: Stages 0 + 1 =====================

% ---- Stage 0: streaming input ----
\node[iobox] (frames) {frames\\@ 1\,fps};
\node[iobox, below=1.5mm of frames] (slam) {SLAM\\$({\rm lat,lng},t)$};

% ---- Stage 1: PSM ingest ----
% Encoder
\node[box, right=5mm of frames] (enc) {VLM encoder\\\textit{CLIP / SigLIP}};
\node[box, right=2mm of enc, fill=blue!8] (emb) {$\mathbf{e}$};

% H3 lookup
\node[box, right=5mm of slam, fill=gray!8] (h3lookup) {H3 lookup\\$r{=}12$};

% H3 hex grid — compact 7-hex cluster
\coordinate (hexcenter) at ($(emb.east) + (12mm, -4mm)$);
\node[hex] (h0) at ($(hexcenter) + (-4.2mm, 2.4mm)$) {};
\node[hex] (h1) at ($(hexcenter) + ( 0mm,   4.8mm)$) {};
\node[hex] (h2) at ($(hexcenter) + ( 4.2mm, 2.4mm)$) {};
\node[hexhi] (hC) at (hexcenter) {};
\node[hex] (h4) at ($(hexcenter) + (-4.2mm,-2.4mm)$) {};
\node[hex] (h5) at ($(hexcenter) + ( 0mm,  -4.8mm)$) {};
\node[hex] (h6) at ($(hexcenter) + ( 4.2mm,-2.4mm)$) {};

% Per-cell state pop-out (right of hex grid)
\coordinate (cellAnchor) at ($(hexcenter) + (14mm, 0mm)$);

% Ring buffer of HLL sketches
\node[hllold]  (hll1) at ($(cellAnchor) + (-3.0mm, 3.5mm)$) {};
\node[hllold]  (hll2) at ($(cellAnchor) + (-1.0mm, 3.5mm)$) {};
\node[hllbox]  (hll3) at ($(cellAnchor) + ( 1.0mm, 3.5mm)$) {};
\node[hllbox]  (hll4) at ($(cellAnchor) + ( 3.0mm, 3.5mm)$) {};
\node[hllbox]  (hll5) at ($(cellAnchor) + ( 5.0mm, 3.5mm)$) {};
\node[font=\tiny, right=0.3mm of hll5, anchor=west, align=left]
     (hlllbl) {HLL ring\\[-0.5ex]\textit{(decay)}};

% Reservoir slots
\node[resvslot] (r1) at ($(cellAnchor) + (-3.0mm, 0mm)$) {};
\node[resvslot] (r2) at ($(cellAnchor) + (-1.0mm, 0mm)$) {};
\node[resvslot] (r3) at ($(cellAnchor) + ( 1.0mm, 0mm)$) {};
\node[resvslot] (r4) at ($(cellAnchor) + ( 3.0mm, 0mm)$) {};
\node[resvslot] (r5) at ($(cellAnchor) + ( 5.0mm, 0mm)$) {};
\node[font=\tiny, right=0.3mm of r5, anchor=west, align=left]
     (rsvlbl) {$R$ exemplars};

% Bounded label
\node[font=\tiny\itshape, below=2mm of r3, gray!60!black]
     (boundlbl) {bounded per cell};

% Connector from highlighted hex to pop-out
\draw[smallflow, dashed, blue!50!black]
     (hC.east) -- ($(hll1.west) + (-0.8mm, -1.8mm)$);

% Stage-0/1 flow arrows
\draw[flow] (frames.east) -- (enc.west);
\draw[flow] (enc.east)    -- (emb.west);
\draw[flow] (slam.east)   -- (h3lookup.west);
\draw[flow] (emb.east) -- ++(2mm,0) |- (hC.west);
\draw[flow] (h3lookup.east) -- ++(2mm,0) |- (h4.west);

% ===================== BOTTOM ROW: Stages 2 + 3 =====================

% Vertical offset below the ingest row
\coordinate (row2anchor) at ($(boundlbl.south) + (0, -6mm)$);

% ---- Stage 2: query ----
\node[iobox, align=center] (query)
    at ($(frames.south) + (0, -20mm)$)
    {Question $q$};
\node[box, right=3mm of query] (qenc) {VLM enc.};
\node[box, right=2mm of qenc, fill=blue!8] (qemb) {$\mathbf{q}$};

\node[box, right=4mm of qemb, align=center, fill=orange!10,
      draw=orange!70!black] (topk)
    {top-$K$ cosine\\{\tiny\texttt{per\_cell\_cap}}};

\node[box, right=3mm of topk, align=center, fill=blue!5] (cands)
    {$K{=}5$ candidates\\$(t,\text{cell},\text{lat},\text{lng})$};

% Flow: reservoir → top-K (vertical connector from row 1)
\draw[flow, dashed, blue!50!black]
    ($(r3.south) + (0, -1mm)$) -- ++(0, -5.5mm) -| (topk.north);

% Stage-2 flow arrows
\draw[flow] (query.east) -- (qenc.west);
\draw[flow] (qenc.east)  -- (qemb.west);
\draw[flow] (qemb.east)  -- (topk.west);
\draw[flow] (topk.east)  -- (cands.west);

% ---- Stage 3: MLLM rerank ----
\node[box, right=4mm of cands, align=center,
      fill=green!8, draw=green!50!black] (fetch)
    {fetch\\frames};
\node[box, right=2mm of fetch, align=center,
      fill=green!12, draw=green!50!black] (mllm)
    {MLLM\\rerank};
\node[iobox, right=2mm of mllm] (ans) {answer};

\draw[flow] (cands.east) -- (fetch.west);
\draw[flow] (fetch.east) -- (mllm.west);
\draw[flow] (mllm.east)  -- (ans.west);

% ===================== Stage bounding boxes =====================
\node[stage, fit=(frames)(slam),
      label={[stagelbl, anchor=south]above:{\tiny Stage 0}}] {};
\node[stage, fit=(enc)(emb)(h3lookup)(h0)(h2)(h5)(rsvlbl)(boundlbl),
      label={[stagelbl, anchor=south]above:{\tiny Stage 1: ingest}}] {};
\node[stage, fit=(query)(qenc)(qemb)(topk)(cands),
      label={[stagelbl, anchor=south]above:{\tiny Stage 2: retrieve}}] {};
\node[stage, fit=(fetch)(mllm)(ans),
      label={[stagelbl, anchor=south]above:{\tiny Stage 3: rerank}}] {};

\end{tikzpicture}
\caption{\textbf{PSM architecture.}
\emph{Top:} streaming frames are encoded and deposited into H3 cells;
each cell keeps a bounded ring of $C$ HLL sketches (orange, fading
with age) and $R$ reservoir-sampled embeddings (blue).
\emph{Bottom:} the query is encoded, cosine-searched against the
reservoirs with a \caponek{} gate, and the top-$K$ candidate frames
are reranked by a frontier MLLM.}
\label{fig:f1-architecture}
\end{figure}

Figure~\ref{fig:f1-architecture} sketches the end-to-end architecture.

\paragraph{Probabilistic Spatial Memory.}
We propose \textbf{\psm{}}, a streaming spatial-temporal index over an H3
hex-grid tessellation with three per-cell structures: reservoir-sampled
embedding exemplars; a time-decayed ring buffer of HyperLogLog sketches giving
per-cell cardinality and a temporal envelope; and a per-cell capacity bound
(\caponek{}) that dials strict spatial diversity ($\texttt{cap}{=}1$) to
unconstrained top-$K$ ($\texttt{cap}{=}K$). Ingestion is streaming and per-cell
state is bounded; query cost depends on scope (a spatially-centered query scans
only a k-ring of cells, while an uncentered query scans all retained exemplars).
We report engine ingest/query cost at the actual CLIP-L deployment
configuration in \S\ref{sec:results-memory-latency}; the CLIP encoder, not the
index, is the dominant on-device cost (\S\ref{sec:limitations}).

\paragraph{Evaluating the retrieval substrate.}
We evaluate \psm{} as a retrieval substrate on egocentric look-back QA (e.g., ``where did I last see my bike?'', ``when was I last at this intersection?''). Multimodal LLMs achieve strong single-frame grounding yet remain brittle on temporal grounding over long video~\cite{cheng2024videgothink,chen2025groundedmultihop}, so we compose \psm{} as an upstream spatial \emph{prefilter} to a frontier MLLM (\S\ref{sec:pipeline}): \psm{} supplies fixed-budget recall, the MLLM supplies localization precision.

\paragraph{Contributions.}
\begin{itemize}
  \item \textbf{A budget-matched allocation study.} On 14 street-scale
    egocentric sessions we compare streaming retention policies at an identical
    global exemplar budget and candidate count (global reservoir, FIFO,
    temporal stratification, a semantic $k$-center coreset, and a causal
    spatial-allocation probe), with rare/common-place and old/recent strata
    reported together (\S\ref{sec:results-main}).
  \item \textbf{A characterization, not a win.} Spatial allocation changes
    \emph{which} places are remembered, not aggregate retrieval: it helps
    rare/low-exposure-place recall in a \emph{place-specific} way (it beats
    coordinate-permutation nulls on 2/3 seeds and is robust to grid
    translation/rotation) but trades against common places, and a non-spatial
    semantic-diversity coreset is best overall. The spatial policy is an
    experimental probe with area-growing bookkeeping, not a deployed method.
  \item \textbf{\psm{} substrate + honest systems accounting.} A streaming H3
    index with per-cell reservoir exemplars and a time-decayed HLL ring
    (\S\ref{sec:method}); we report the engine's on-device cost at the actual
    CLIP-L deployment configuration and the corrected per-cell memory model on
    102.9\,h of driving (Honda HDD) as modeled logical state
    (\S\ref{sec:results-hdd}).
  \item \textbf{Correcting the MLLM framing.} A frontier-MLLM comparison shows
    query-aware retrieval beats uniform frame sampling; it does \emph{not} show
    that MLLMs require \psm{} (\S\ref{sec:results-vanilla-mllm}). Engine and
    budget-matched harness (policies, semantic/coordinate-null controls, paired
    statistics) are open-source under MIT.
\end{itemize}
